\documentclass[11pt,a4paper]{article}

\usepackage[utf8]{inputenc}

\usepackage[T1]{fontenc}

\usepackage[polish]{babel}

\usepackage{lmodern}

\usepackage{geometry}

\geometry{margin=27mm}

\usepackage[hidelinks]{hyperref}

\title{Pamięć prospektywna po przerwaniu zadania: przypomnienia kontekstowe wobec czasowych}

\author{Mateusz Pomianek}

\date{Szkic koncepcyjny, wersja 0.1.2}

\begin{document}

\maketitle

\begin{abstract}

Tematem pracy jest Pamięć prospektywna po przerwaniu zadania: przypomnienia kontekstowe wobec czasowych. Tekst systematyzuje uzasadnienie, hipotezy i sposób weryfikacji zagadnienia. Proponowana miara trafności realizacji odroczonych intencji wymaga formalnego zatwierdzenia w protokole. Koncepcja pozostaje w stanie \texttt{NEW\_CROSS\_PROGRAM\_IDEA\_HOLD}; nie zebrano danych ani nie przedstawiono wyników.

\end{abstract}

\section{Status i granice własności}

Właścicielem naukowym i autorem niniejszego szkicu jest Mateusz Pomianek. Identyfikator koncepcji to \texttt{MSC-AAL-P08}; praca pozostaje propozycją badawczą wstrzymaną na etapie \texttt{IDEA/HOLD}. Miejsce przechowywania szkicu w programie ARC Neuro nie przesądza o kanonicznym repozytorium przyszłego protokołu i danych. Przypisanie takiego repozytorium wymaga odrębnego potwierdzenia w rejestrze portfolio.

\section{Wprowadzenie}

Samodzielność osądu i funkcje poznawcze zależą od jakości informacji oraz sposobu jej przedstawienia. Zagadnienie Pamięć prospektywna po przerwaniu zadania: przypomnienia kontekstowe wobec czasowych wymaga odróżnienia badanego konstruktu od jego operacjonalizacji i kontekstowych źródeł zmienności. Trzeba porównać możliwe wyjaśnienia i ustalić granice interpretacji dowodów.

\section{Cel i pytania badawcze}

Celem pracy jest rozstrzygnięcie zagadnienia: Pamięć prospektywna po przerwaniu zadania: przypomnienia kontekstowe wobec czasowych. Dobór pierwszorzędowego wskaźnika pozostaje elementem przyszłego protokołu. Hipotezy i ich kierunki muszą zostać określone przed analizą danych lub krytyczną syntezą źródeł.

\section{Materiał i metody}

W scenariuszu domowym 2D porównać przypomnienia kontekstowe i czasowe przy równej liczbie komunikatów. Jako kandydat na pierwszorzędowy wskaźnik rozważana jest proporcja poprawnie wykonanych odroczonych intencji w ustalonym oknie. Definicja okna, jednostki analizy i reguł uznania odpowiedzi za poprawną wymaga zatwierdzenia przed zbieraniem danych. Nie zebrano danych w ramach tego szkicu. Przed rozpoczęciem badania z udziałem ludzi należy ustalić plan doboru i liczebności próby, wymogi etyczne, procedury zgody i ochrony danych.

\section{Plan analizy}

Proponowany plan analityczny: regresja logistyczna z efektem uczestnika i zadania oraz kontrolą przerwań. Zamrożenie estymandu, kryteriów jakości i analiz wrażliwości nastąpi przed dostępem do danych; na tym etapie nie ma oszacowań ani wyników.

\section{Warunki wznowienia koncepcji}

Przed zmianą statusu \texttt{HOLD} należy ustalić właściciela kanonicznego projektu, rozgraniczyć temat względem powiązanych koncepcji, uzasadnić nowość na podstawie zweryfikowanych źródeł, uzgodnić wskaźnik główny i estymand oraz przygotować protokół, plan liczebności próby, analiz i ochrony danych. Dostępność stanowiska oraz oprogramowania pozostaje niepotwierdzona.

\section{Wyniki}

W niniejszym pliku nie umieszczono wyników empirycznych, statystyk ani nieprzeprowadzonych analiz. Sekcja wymaga uzupełnienia po uzyskaniu i weryfikacji materiału dowodowego.

\section{Dyskusja i ograniczenia}

Ograniczenia obejmują możliwe błędy pomiaru, dobór próby, alternatywne mechanizmy i ograniczoną możliwość uogólniania. Korelacja nie dowodzi przyczynowości; poprawność techniczna nie jest dowodem znaczenia klinicznego.

\section{Wnioski}

Zdefiniowano ramy opracowania, ale nie sformułowano wniosków wykraczających poza dostępny stan dowodów.

\section*{Bibliografia}

Weryfikacja literatury i przypisanie cytowań do stwierdzeń są kolejnym etapem prac.

\end{document}

